\documentclass[11pt]{article}

\usepackage[margin=1in]{geometry}
\usepackage{amsmath,amssymb}
\usepackage{booktabs}
\usepackage{graphicx}

\usepackage{microtype}
\usepackage[colorlinks=true,linkcolor=blue,citecolor=blue,urlcolor=blue]{hyperref}

\newcommand{\IG}{\mathrm{IG}}
\newcommand{\Ent}{H}

\title{Planning the Next Evidence Acquisition Step:\\
An Offline Framework for Active Scientific Discovery}

\author{Hui Mao \\
\normalsize Independent Researcher \\
\normalsize \texttt{hui.mao@alumni.upenn.edu}}
\date{August 2026}

\begin{document}

\maketitle
\thispagestyle{empty}
\renewcommand{\thefootnote}{\fnsymbol{footnote}}
\footnotetext[1]{Fourth in a series on temporally grounded evaluation of
scientific question discovery; see~\cite{paper1},~\cite{paper2},
and~\cite{paper3}. Code, configurations, and the scaled benchmark are
public (Reproducibility).}
\renewcommand{\thefootnote}{\arabic{footnote}}

\begin{abstract}
Most work on machine-assisted science focuses on generating scientific
questions or executing experiments autonomously. We address the step in
between: given an existing scientific question, a set of competing
hypotheses, and the evidence currently available, \emph{what should we
investigate next to maximize expected scientific progress?} We present an
offline, domain-independent planning framework that represents evidence as
a typed, confidence-weighted graph; identifies hypotheses with missing,
weak, or non-independent evidence; generates generic candidate
evidence-acquisition actions; and ranks them by a configurable utility
combining expected information gain, novelty, tractability, and cost.
Beliefs over hypotheses are revised with a modular Bayesian updater,
closing the plan--observe--update loop without ever executing actions.
Because the planner is offline, it can be evaluated honestly against
history: a benchmark truncates the evidence graph at a time cutoff, hides
later evidence as ground truth, and asks whether the recommended actions
would have resolved the question earlier. On a worked case study from
exoplanet atmospheric characterization --- a reported strongly subsolar
water abundance whose premise was later refuted by three independent
re-analyses --- the planner's top recommendation is exactly the
discriminating re-analysis that historically resolved the question
(resolution efficiency $1.0$; replaying the refuting evidence removes
$1.06$ bits of hypothesis entropy). We further ground the utility
components in the historical-outcome findings of the companion
predictor-discovery study (Paper 3, $n{=}980$ questions at five
cutoffs): explicitly articulated competing hypotheses predict
resolution, and the resolution channel runs through quantified,
explicitly contradictory, recurrent evidence rather than methodological
doubt --- findings we convert into an empirically calibrated weight
configuration that sharpens the planner's ranking toward the
historically resolving actions. Scaling the benchmark to all $374$ of
Paper 3's tension questions --- each evidence cluster batch-converted
to a graph by transparent rules and planned at its own historical
cutoff --- turns the metrics into distributions and surfaces a
structural result: planner utility runs opposite to measured community
attention (pooled permutation $p = 0.004$), because both are driven,
with opposite sign, by the breadth of the existing evidence base. We
then dismantle the appealing reading of that result. The effect does
not survive stratification by breadth, and its residual concentrates in
the high-specificity stratum where Paper 3's retrieval-bounded label is
documented to under-count engagement. Re-adjudicating all $62$
high-specificity questions blind, against citation-graph retrieval
rather than abstract keyword matching, recovers $94\%$ of the questions
the label calls unaddressed ($69\%$ under the strictest bar) and the
utility effect vanishes ($p = 0.70$ and $0.61$, against $0.004$). The
divergence between resolution-driven and attention-driven allocation is
therefore real as a mechanism but is not evidence of detected neglect;
the durable empirical result is a measured bound on the outcome label
itself. The framework is released as a reference implementation with
seven evaluation metrics, full configurability, and no domain
assumptions in code.
\end{abstract}

\section{Introduction}

Scientific progress is rarely limited by the ability to ask questions; it
is limited by the cost of acquiring the evidence that discriminates
between competing answers. A researcher facing an open question --- is a
reported effect real, an artifact of the analysis, a confounder, or an
instrument systematic? --- must decide where to spend the next unit of
effort: re-read the literature, re-run an analysis with different
assumptions, compare independent datasets, run a simulation, or invest in
a costly new observation. This decision is made under uncertainty, under
resource constraints, and usually informally.

This paper treats that decision as a planning problem. We assume the
scientific question already exists (produced by a human or by an upstream
question-generation system) and ask only: \emph{given the current
evidence, which evidence-acquisition step maximizes expected scientific
progress?}

This is Paper 4 of a series. Paper 1 built an evidence-graph system that
generates questions from cross-paper observational tensions; Paper 2
introduced historical backtesting of question generators (freeze
questions on pre-cutoff literature, observe what the community actually
did afterwards); Paper 3\footnote{\url{https://github.com/nonameisready/scientific-question-outcome-prediction-indicators}}
scaled that protocol into a supervised predictor-discovery study over
$980$ astronomy questions frozen at five cutoffs (2012--2020) on the
complete arXiv astro-ph record, with strict blinded outcome labels.
Those papers ask which \emph{questions} the future rewards; this paper
takes the question as given and plans the \emph{evidence acquisition}
that follows. Paper 3's outcome dataset feeds back into this work in two
concrete ways developed in Section~\ref{sec:grounding}: its controlled
effects empirically ground (and partly correct) our utility components,
and its released questions and outcome labels are directly importable as
benchmark units for our evaluation protocol.

Our design commitments differ from much recent work on autonomous
science:

\begin{itemize}
  \item \textbf{Offline.} The system is a planner, not an agent. It does
  not browse, does not call instrument APIs, and does not execute
  experiments. Everything works from static datasets, which makes the
  system auditable and makes rigorous historical evaluation possible.
  \item \textbf{Domain-independent.} No scientific field is hardcoded.
  Domain knowledge enters only through data (questions, hypotheses,
  evidence graphs) and configuration (action catalogs, weights).
  \item \textbf{Modular.} Every component --- hypothesis generation, gap
  analysis, action generation, each utility term, belief updating,
  evaluation --- sits behind a small interface and can be replaced
  independently.
\end{itemize}

\paragraph{Contributions.}
(1)~A formalization of next-step evidence planning over typed,
confidence-weighted evidence graphs with provenance;
(2)~a gap analysis that flags hypotheses with missing, sparse, weak, or
non-independent evidence;
(3)~a configurable utility decomposition --- expected information gain via
preposterior analysis, provenance-based novelty, resource-aware
tractability, and cost;
(4)~a historical-simulation benchmark with seven metrics that evaluates
whether a planner would have resolved a question earlier;
(5)~an empirical grounding of the utility components in Paper 3's
historical-outcome dataset, including an import path for its released
questions and labels and an empirically calibrated weight
configuration;
(6)~a scaled benchmark run over all $374$ of Paper 3's tension
questions, batch-converted to evidence graphs by auditable rules, with
metric distributions and an outcome-stratified analysis;
(7)~a blind re-adjudication of the $62$ questions carrying that
analysis' residual effect, which measures the outcome label's
under-counting at $69$--$94\%$, withdraws our own neglect
interpretation, and bounds what specificity-related findings computed
on that label can mean; and
(8)~an open reference implementation with unit tests and a worked case
study whose hidden ground truth follows the documented historical
outcome.

\section{Related work}

\paragraph{Choosing experiments.} The problem of allocating the next
increment of scientific effort has a direct precedent in the science of
science. Rzhetsky et al.~\cite{rzhetsky2015} model the choice of
experiments as a search over a knowledge network and show that the
strategies scientists actually adopt are collectively suboptimal ---
too concentrated near established results --- and that a planner
optimizing collective discovery would allocate differently. Foster et
al.~\cite{foster2015} document the same conservatism at the level of
research strategy. Our framework sits at the operational end of this
line: rather than characterizing aggregate strategy, it plans one
concrete next action for one question against an explicit hypothesis
set, and the divergence we measure in Section~\ref{sec:scale} between
planner utility and realized community attention is a question-level
instance of exactly the gap those papers describe at population level.

\paragraph{Autonomous science.} A parallel line builds systems that
execute research end to end, from the Nobel Turing
Challenge~\cite{kitano2021} onward. We deliberately stop short of
execution. The planner is an offline decision layer that can precede
any executor, human or automated, which is what makes its
recommendations auditable and its evaluation possible against history
rather than against a simulator.

\paragraph{Evaluating generated ideas.} Novelty-based evaluation of
machine-generated research ideas has repeatedly failed to predict value.
Si et al.~\cite{si2024novel} measured expert opinion at generation
time; their follow-up~\cite{si2025execution} had researchers execute
both LLM and human ideas and found the ideas rated highest before
execution scored lowest after. HindSight~\cite{hindsight} matches
generated ideas against subsequent literature and finds LLM-judged
novelty \emph{negatively} correlated with realized impact. Papers 1--3
of this series~\cite{paper1,paper2,paper3} replace opinion with
history at the level of questions; this paper carries the same
substitution into the planning step, scoring plans by whether the
evidence they would acquire resolves anything.

\section{Problem Statement}

Let $q$ be a scientific question and
$\mathcal{H} = \{h_1, \dots, h_n\}$ a set of mutually exclusive,
exhaustive competing hypotheses with probabilities
$p = (p_1, \dots, p_n)$, $\sum_i p_i = 1$. Let $G = (V, E)$ be the
current \emph{evidence graph} (Section~\ref{sec:graph}). Let
$\mathcal{A}$ be a set of candidate evidence-acquisition actions, each
with a type, relative cost, relative time, required resources, and a
target hypothesis.

The planner must return a ranking of $\mathcal{A}$ (top-$K$
recommendations) maximizing a utility
\begin{equation}
U(a) \;=\; w_{\mathrm{ig}}\,\IG(a) \;+\; w_{\mathrm{nov}}\,\mathrm{Nov}(a)
\;+\; w_{\mathrm{tr}}\,\mathrm{Tr}(a) \;-\; w_{c}\,C(a),
\label{eq:utility}
\end{equation}
where $\IG$ is the expected reduction in hypothesis uncertainty,
$\mathrm{Nov}$ measures independence from existing evidence,
$\mathrm{Tr}$ measures feasibility, $C$ is cost, and the weights $w$ are
configurable. After an action is (externally) executed and its outcome
observed, hypothesis probabilities are revised and the cycle repeats. The
framework plans and updates; it never executes.

\section{Framework}

\subsection{Evidence Graphs and Provenance}
\label{sec:graph}

Evidence is a directed multigraph with five node types ---
\emph{paper}, \emph{observation}, \emph{dataset}, \emph{claim},
\emph{hypothesis} --- and four edge types --- \emph{supports},
\emph{contradicts}, \emph{derived\_from}, \emph{observes} --- each edge
carrying a confidence in $[0,1]$. Claims point at the hypotheses they
support or contradict; derived nodes point at their origins
(claim $\to$ paper $\to$ dataset); observations point at what they
observe. Nodes optionally carry timestamps, which enables historical
simulation (Section~\ref{sec:eval}).

Two queries are load-bearing. First,
$\mathrm{evidence}(h)$ returns the supports/contradicts edges incident to
a hypothesis. Second, $\mathrm{roots}(v)$ follows
\emph{derived\_from} edges transitively to the root sources of a node.
Two claims that descend from the same dataset are \emph{not} independent
evidence, no matter how many papers separate them; this notion of
provenance drives both gap analysis and novelty.

\subsection{Hypothesis Generation}

Hypotheses are structured objects (identifier, description, prior,
supporting and contradicting evidence references). The reference
implementation provides two interchangeable generators behind one
interface: a \emph{static} generator that loads curated hypothesis sets
keyed by question (the normal path for imported datasets), and a
\emph{template} generator producing four generic methodological
alternatives applicable to any empirical claim: the effect is real; it is
an artifact of the analysis method; it is confounded by an unmodeled
factor; it is an instrument or measurement systematic. Priors are
normalized, defaulting to uniform.

\subsection{Evidence Gap Analysis}
\label{sec:gaps}

For each hypothesis the analyzer emits the first applicable gap:

\begin{enumerate}
  \item \textbf{No evidence.} No edge touches the hypothesis
  (gap confidence $1.0$).
  \item \textbf{Sparse evidence.} Fewer than $k_{\min}$ evidence edges;
  confidence $0.5 + 0.5\,(1 - \mathrm{count}/k_{\min})$.
  \item \textbf{Weak support.} Summed support confidence below a
  threshold $s_{\min}$; confidence $0.4 + 0.4\,\delta$ where $\delta$ is
  the relative deficit.
  \item \textbf{No independent corroboration.} All evidence descends from
  a single provenance root (confidence $0.78$): formally,
  $\big|\bigcup_{v \in \mathrm{sources}(h)} \mathrm{roots}(v)\big| \le 1$.
\end{enumerate}

The output is a gap report: hypothesis, missing-evidence description,
confidence, and a human-readable reason, sorted most-confident first. All
thresholds are configurable.

\subsection{Candidate Action Generation}

Actions are instantiated from a generic, YAML-configurable catalog of
seven types: read paper, search archive, re-run analysis, collect
observation, run simulation, compare datasets, search review. Each
template carries a relative cost and time in $[0,1]$ and required
resource tags. For each gap, the generator instantiates the subset of
types that plausibly closes that kind of gap --- independence gaps
receive independent-source actions (collect observation, compare
datasets, search archive); weak-support gaps receive re-analysis actions;
evidence-free hypotheses receive broad search actions --- and records the
gap's reason and confidence in the action's metadata. Nothing in the
catalog or the mapping references any scientific domain.

\subsection{Utility Estimation}
\label{sec:utility}

Each term of Eq.~\eqref{eq:utility} is an independent estimator behind a
common interface, injected into the combiner; any term can be replaced
without touching the others.

\paragraph{Expected information gain (preposterior analysis).}
An action $a$ targeting hypothesis $h_t$ is modeled as producing a binary
outcome $o \in \{\text{favors } h_t,\ \text{disfavors } h_t\}$ with a
likelihood ratio $r_a$ determined by the action type's
\emph{discriminative power} (a configurable table; an independent
observation separates hypotheses more sharply than reading a paper):
\begin{equation}
P(o^+ \mid h_t) = \frac{r_a}{r_a + 1}, \qquad
P(o^+ \mid h_j) = \frac{1}{r_a + 1} \quad (j \ne t).
\end{equation}
The estimate is the expected entropy reduction
\begin{equation}
\IG(a) \;=\; \Ent(p) \;-\; \mathbb{E}_{o}\!\left[\Ent\!\big(p \mid o\big)\right]
\;\ge\; 0,
\end{equation}
with $\Ent$ the Shannon entropy in bits and the posterior computed by
Bayes' rule for each outcome. This is exact for the two-outcome model and
degrades gracefully: actions with likelihood ratio $1$ yield zero gain,
as do actions with no target among the live hypotheses.

\paragraph{Novelty (provenance independence).}
Each action type has a base source novelty $\nu_a \in [0,1]$ (new
observation $=1.0$; re-reading literature $\approx 0.3$). The base is
boosted by the provenance concentration of the target's existing
evidence:
\begin{equation}
\mathrm{Nov}(a) \;=\; \min\!\Big(\nu_a \big(1 + \tfrac{1}{\max(|R_t|,\,1)}\big),\; 1\Big),
\end{equation}
where $R_t$ is the set of provenance roots behind the target hypothesis's
evidence. If everything known descends from one dataset, independent
evidence is maximally valuable.

\paragraph{Tractability.}
$\mathrm{Tr}(a) = (1 - \mathrm{time}_a) \cdot \rho^{m}$, where $m$ is the
number of required resources not declared available and $\rho$ (default
$0.5$) is the missing-resource penalty. With no declared resource set,
all resources are assumed available.

\paragraph{Cost.}
The action's configured relative cost, passed through.

\paragraph{Calibration.}
Information gain is measured in bits and is typically well below $1$ for
a single action, while the other terms are $[0,1]$ scores; the default
configuration therefore sets $w_{\mathrm{ig}} = 10$ to place the terms on
a comparable scale, with all weights exposed in YAML.

\subsection{Planning}

The planner is deliberately thin: it scores every candidate with the
injected utility estimator, sorts by total utility, truncates to
top-$K$, and attaches to each recommendation its utility breakdown, its
target hypothesis, and a reasoning string naming the dominant term. All
ranking policy lives in the (replaceable) utility estimator.

\subsection{Belief Updating}

New evidence arrives as a likelihood vector
$L = \big(P(e \mid h_1), \dots, P(e \mid h_n)\big)$; hypotheses absent
from the vector receive a neutral likelihood of $1$. The default updater
is exact discrete Bayes,
\begin{equation}
p_i' \;=\; \frac{p_i\, P(e \mid h_i)}{\sum_j p_j\, P(e \mid h_j)},
\end{equation}
with two safeguards: evidence impossible under every hypothesis leaves
the distribution unchanged (it is uninformative for this hypothesis
set), and the update is non-destructive, recording the evidence
identifier on each hypothesis's supporting or contradicting list
according to the direction of the probability change. The updater is a
protocol; richer models (soft evidence, correlated hypotheses,
hierarchical priors) can be substituted without touching callers.

\section{Empirical Grounding from Historical Outcomes (Paper 3)}
\label{sec:grounding}

The utility decomposition of Section~\ref{sec:utility} was derived from
first principles. Paper 3's predictor-discovery study now lets us check
those principles against a thousand historical outcomes: $980$ questions
frozen at five cutoffs, each labeled by a strict blinded judge with what
the following five years of literature actually did (52.7\% substantively
addressed; interpretable models predict 2020 outcomes at AUC $0.71$
out-of-time and out-of-domain). Four of its findings bear directly on
this framework.

\paragraph{Explicit competing hypotheses are the one falsifiability
property with predictive teeth.} Across controlled analyses, explicitly
articulated competing hypotheses raise the probability of future
attention ($\beta = +0.17$, $p = 0.02$) and, directionally, of being
\emph{answered} given attention ($\beta = 0.56$, $p = 0.08$) --- the only
falsifiability feature that survives Paper 3's four-way robustness
screen. This is direct empirical support for the hypothesis-decomposition
step this framework makes mandatory: a question is planned \emph{as} a
set of competing hypotheses, never as free text.

\paragraph{The resolution channel is discriminative, not diffuse.}
Within Paper 3's tension-mined subset (374 questions with verbatim
evidence-graph provenance), every premise refutation comes from a
\emph{quantified} tension cluster ($3/63$ quantified vs.\ $0/311$
unquantified), five of six resolutions are explicit-contradiction-typed,
all six resolutions come from tensions \emph{recurrent} across cutoffs
--- and clusters built on methodological-challenge language engaged the
community at the average rate yet resolved \emph{nothing} ($0/61$).
Translated into this framework's terms: evidence with a high likelihood
ratio against a specific, contestable claim is what closes questions;
diffuse methodological worry produces attention without resolution. This
is precisely the behavior encoded in the information-gain estimator's
discriminative-power table, and the reason our utility centers expected
entropy reduction rather than engagement.

\paragraph{Attention is community-relational; novelty is weak and
redundant.} Paper 3's robust predictors of future attention are prior
community recognition (overlap with pre-cutoff reviews,
$\beta = +0.34$), the density of directly related prior evidence
($\beta = +0.26$), and cluster breadth within the tension subset
($\beta = +0.45$/SD). Its verdicts on novelty and tension are more
specific than the slogan they are often reduced to, and worth stating
exactly: semantic novelty carries a \emph{small positive} controlled
effect ($\beta = +0.17$, $p = 0.03$) that is only half sign-stable
across generation sources, and the novelty feature group is redundant
given the rest (dropping it \emph{improves} held-out AUC, $0.688 \to
0.708$); what is null is evidence \emph{tension density} itself
($\beta = 0.06$, $p = 0.52$). Five of six pre-stated mechanism
hypotheses failed, including ``evidence tension beats novelty'' ---
which failed because tension was null, not because novelty was.
Independently, HindSight~\cite{hindsight} finds LLM-judged novelty
\emph{negatively} correlated with realised impact, and the
ideation--execution study of Si et al.~\cite{si2025execution} finds
LLM ideas rated highest before execution scoring lowest after: three
different designs converge on novelty being a poor optimization target
even where it is not a null one.

For this framework the lesson is a separation of objectives. Our
planner optimizes expected \emph{resolution} (entropy reduction), not
expected attention, and its novelty term is not semantic distance at
all but \emph{provenance independence} --- the step from one evidence
root to a second. The closest analogue to that construct in Paper 3's
feature set is instrument diversity of the existing evidence base,
whose controlled effect on attention is \emph{negative}
($\beta = -0.21$): multi-instrument questions read as synthesis-level
problems no single group owns. That is the honest empirical basis for
down-weighting our novelty term --- not a claim that novelty predicts
nothing --- while its role in the refutation channel (quantified
cross-checks of a single-source claim) is what keeps the term in the
utility at all. We accordingly ship an empirically calibrated
configuration
(\texttt{configs/paper3\_informed.yaml}) that keeps the
information-gain weight, halves the novelty weight, and raises
tractability (Paper 3's Cox model: archival-data availability
accelerates engagement; requiring a not-yet-operational instrument is
the largest negative effect). Section~\ref{sec:case} shows its effect
on the case-study ranking.

\paragraph{A shared benchmark substrate.} Paper 3 releases its
questions, provenance clusters, and tiered outcome labels. Its records
load directly through our question loader (an importer ships in
\texttt{experiments/}), and each labeled question is a candidate
benchmark unit for Section~\ref{sec:eval}: the cutoff supplies the
historical-simulation boundary and the label supplies hidden ground
truth. Its learned structure prior
(\texttt{results/structure\_prior.json}, structure-only AUC $0.586$
vs.\ $0.502$ for the full question-level record within a single source)
is likewise importable as an empirical prior over which evidence
configurations the community will engage --- a community-uptake signal
complementary to our resolution-centric utility.

\section{Offline Evaluation by Historical Simulation}
\label{sec:eval}

An offline planner can be evaluated the only way that avoids
self-grading: against history. The benchmark takes a \emph{full}
timestamped evidence graph and a cutoff $\tau$:

\begin{enumerate}
  \item Filter the graph to nodes with timestamp $< \tau$
  (untimestamped structural nodes are kept; edges touching hidden nodes
  are dropped). The planner sees only this sub-graph.
  \item Run any planner pipeline (a two-method protocol) on the visible
  sub-graph, producing recommendations and the hypothesis set used.
  \item Express post-cutoff evidence as ground-truth records: which
  hypothesis it bore on, which action type historically produced it, and
  its per-hypothesis likelihoods.
  \item A ground-truth item is \emph{reached} if some recommendation
  matches its action type and target hypothesis; reached items are
  replayed through the belief updater to measure the uncertainty the plan
  could have removed.
\end{enumerate}

The central question is: \emph{did the planner recommend the actions that
would have resolved the question earlier?}

\subsection{Metrics}

\begin{description}
  \item[Top-$K$ utility.] Mean total utility of the top-$K$
  recommendations.
  \item[Hypothesis entropy reduction.] $\Ent(p_{\text{before}}) -
  \Ent(p_{\text{after}})$ after replaying reached ground truth ---
  what history says the plan was worth.
  \item[Average information gain.] Mean \emph{expected} gain across
  recommendations --- what the planner believed; comparing it with
  realized entropy reduction exposes over- or under-confidence.
  \item[Action diversity.] Distinct action types over number of
  recommendations, penalizing degenerate plans.
  \item[Coverage of competing hypotheses.] Fraction of hypotheses
  targeted by at least one recommendation; a good plan probes the space
  of explanations, not just the front-runner.
  \item[Resolution efficiency.] Fraction of ground-truth evidence
  reached by the plan; $1.0$ means every piece of evidence that
  historically resolved the question was on the recommended path.
  \item[Expected utility.] Rank-discounted utility with weights
  $1/\mathrm{rank}$ (normalized), modeling that earlier recommendations
  are likelier to be executed.
\end{description}

\section{Case Study}
\label{sec:case}

We instantiate the framework on a question from exoplanet atmospheric
characterization --- \emph{``Is the reported low water abundance real or
a retrieval artifact?''} --- stressing that the domain lives entirely in
the data files; the code is untouched. The question is not arbitrary: it
is the question family of Paper 2's pilot benchmark, whose documented
history (carried forward into Paper 3's dataset) records that the
reported strongly subsolar water abundance was a premise \emph{later
refuted by three independent re-analyses} --- exactly the convergence the
question called for. Our hidden ground truth follows that history: three
post-cutoff evidence items (two independent retrieval re-analyses and
one cross-dataset comparison, 2020--2021) that support the
retrieval-artifact hypothesis and contradict the ``real'' one.

Four hypotheses compete: the low abundance is real ($p=0.3$); it is a
retrieval artifact ($p=0.3$); it is a cloud--composition degeneracy
($p=0.2$); it is an instrument systematic ($p=0.2$). The pre-cutoff
graph (cutoff $\tau = 2020$; the planner sees 10 of 16 nodes) contains
the discovery paper and an early re-analysis, both derived from the
\emph{same} transmission spectrum dataset --- textbook provenance
concentration.

Gap analysis flags the ``real'' hypothesis for lacking independent
corroboration (all evidence descends from one dataset), the ``artifact''
hypothesis for sparse evidence, and the remaining two for having no
evidence. Under the default configuration the planner's top
recommendation is re-running the retrieval under alternative settings,
targeting the artifact hypothesis --- the action family that historically
refuted the premise --- followed by independent observation and
cross-dataset comparison for the two leading hypotheses.

\begin{table}[t]
\centering
\caption{Benchmark results for the case study (top-$K{=}5$, cutoff 2020,
default configuration).}
\label{tab:results}
\begin{tabular}{lr}
\toprule
Metric & Value \\
\midrule
Top-$K$ utility & 2.596 \\
Hypothesis entropy reduction (bits) & 1.062 \\
Average information gain (bits) & 0.173 \\
Action diversity & 0.600 \\
Coverage of competing hypotheses & 0.500 \\
Resolution efficiency & 1.000 \\
Expected utility & 2.616 \\
\bottomrule
\end{tabular}
\end{table}

Table~\ref{tab:results} summarizes the run. Resolution efficiency of
$1.0$ indicates every piece of historically decisive evidence was on the
recommended path; replaying the three refuting items shifts the belief
mass decisively onto the artifact hypothesis, removing $1.06$ of the
$1.97$ bits of initial hypothesis entropy. Coverage of $0.5$ reflects
that the two low-prior hypotheses did not make the top-5 under default
weights --- raising the coverage-relevant weights or $K$ trades focus for
breadth, and the benchmark makes that trade-off measurable.

Re-ranking the same candidates under the Paper 3--informed configuration
of Section~\ref{sec:grounding} sharpens the match with history further:
the discriminating re-analysis stays first, cross-dataset comparison
rises to second and third, and the expensive new observation --- which
history shows was \emph{not} the resolving step for this question ---
drops to the bottom of the top-5. An empirical prior learned from a
thousand historical outcomes and a utility derived from first principles
agree on the head of the ranking and disagree exactly where Paper 3's
evidence says they should: on how much raw source novelty is worth.

The utility weights also behave as designed: with all weights at $1$,
cheap archive searches dominate the ranking and the decisive (expensive)
observation is never recommended; the calibrated default
($w_{\mathrm{ig}}=10$) restores the balance between bits and $[0,1]$
scores. We report this deliberately --- the failure mode is exactly the
kind of miscalibration the benchmark exists to expose.

\section{Scaling Up: 374 Historical Tension Questions}
\label{sec:scale}

The case study is one fully specified benchmark unit. To turn the
metrics into distributions, we batch-convert every evidence-tension
question in Paper 3's released dataset --- 374 questions across eight
astronomy subfields and five cutoffs, each storing its generating
evidence cluster verbatim (sentences, arXiv ids, publication dates;
244 four-paper, 59 three-paper, 71 two-paper clusters) together with a
blinded future-outcome label and cluster-level structure flags.

\paragraph{Conversion.} Each cluster becomes an evidence graph by
transparent, auditable rules (implemented in
\texttt{experiments/paper3\_benchmark.py}). The four template
hypotheses receive priors seeded from the cluster's structure flags:
quantified tensions (scaled by their $N$-sigma) shift mass toward
``real effect'', contradiction and method-challenge language toward
``analysis artifact'', detection-vs-limit stance opposition toward
``unmodeled confounder'', and cross-facility clusters toward
``measurement systematic''. Every cluster sentence becomes a claim node
derived from its (timestamped) paper node, with support edges assigned
by the same flags; hypotheses whose explanatory channel has no flag
remain evidence-free, which is exactly what gap analysis exists to
surface. Distinct arXiv ids count as distinct provenance roots --- an
upper bound on independence, since shared underlying datasets are not
visible at this granularity.

\paragraph{Results.} The planner runs the full pipeline at each
question's own historical cutoff. Table~\ref{tab:scale} reports the
distribution of every at-scale-computable metric. Resolution efficiency
is deliberately absent: Paper 3's labels record \emph{whether and when}
a question was addressed, not \emph{through which action type}, so the
matching that metric requires is not derivable at scale --- mining
action types from the resolving papers is the concrete next step. For
the premise-refuted cases the label does fix the direction of the
hidden evidence; all three trace to a single physical case (the FRB
Galactic-latitude discrepancy, independently re-mined at three
successive cutoffs), for which the direction-only replay removes
$0.095$ bits each.

\begin{table}[t]
\centering
\caption{Metric distributions over the 374 tension questions
(top-$K{=}5$, default configuration, each question at its own cutoff).}
\label{tab:scale}
\begin{tabular}{lrrrr}
\toprule
Metric & Mean & Std & Median & IQR \\
\midrule
Top-$K$ utility & 1.904 & 0.146 & 1.923 & [1.902, 1.923] \\
Average information gain (bits) & 0.056 & 0.007 & 0.056 & [0.054, 0.056] \\
Action diversity & 0.604 & 0.131 & 0.600 & [0.600, 0.600] \\
Coverage of competing hypotheses & 0.505 & 0.110 & 0.500 & [0.500, 0.500] \\
Expected utility & 1.932 & 0.147 & 1.952 & [1.925, 1.952] \\
\bottomrule
\end{tabular}
\end{table}

\paragraph{Planner utility runs opposite to measured attention.}
Stratifying by Paper 3's outcome label, questions the community
subsequently \emph{addressed} score systematically \emph{lower} planner
utility than questions it did not (top-$K$ utility $1.887$ vs.\ $1.919$,
seeded permutation test $p = 0.004$; expected utility $1.916$ vs.\
$1.946$, $p = 0.004$; hypothesis coverage $0.488$ vs.\ $0.520$,
$p = 0.007$). The mechanism is visible in the conversion: broader,
richer evidence clusters yield denser graphs, fewer under-evidenced
hypotheses, and therefore fewer high-utility open actions
(utility vs.\ cluster breadth, Spearman $\rho = -0.24$), while breadth
is simultaneously Paper 3's strongest structural predictor of
\emph{attention} ($\beta = +0.45$/SD). The planner and the community
are, mechanically, reading opposite signs off the same feature.

\paragraph{The effect is breadth-mediated and specificity-confined.}
The tempting reading --- that the planner detects genuine neglect ---
does not survive the robustness checks that Paper 3's own limitations
demand. Two decompositions (Table~\ref{tab:robust}) undercut it. First,
the pooled effect is largely \emph{breadth-mediated}: comparing
addressed against unaddressed questions \emph{within} strata of equal
cluster breadth, the difference is not significant in any stratum
($p = 0.25$, $0.88$, $0.08$). Second, and more sharply, the residual
effect is \emph{concentrated exactly where Paper 3 says its label is
least trustworthy}: within high-specificity questions the gap is large
and significant ($1.830$ vs.\ $1.925$, $p = 0.004$), while in
low-specificity questions it is small and null ($p = 0.11$). Paper 3
documents that its abstract-level retrieval under-counts engagement
non-uniformly, precisely for narrow, object-level, entity-dense
questions --- its own Paper 1 import scores $2/10$ addressed under this
retrieval versus $10/10$ under full-text adjudication, and re-judging
those ten questions on a $2.4\times$ larger corpus flipped four labels
in both directions. A planner that scores evidence-poor, specific
questions highly will therefore anti-correlate with a label that
under-measures exactly those questions, with or without any real
neglect. Rather than caveat this, we measured it.

\subsection{Adjudicating the confounded stratum}
\label{sec:adjudication}

We re-adjudicated all $62$ high-specificity questions against a
retrieval channel the original label did not have. Paper 3 retrieves
candidates by TF--IDF over subfield abstracts, keeping the top eight;
we instead retrieve through the \emph{citation graph}, collecting every
paper that cites a question's generating evidence within its own
five-year post-cutoff window (Semantic Scholar; median $51$ candidates
per question, 4{,}652 in total). Work that engages a question grounded in specific papers
almost always cites one of them, and that path is invisible to keyword
matching on abstracts. Verdicts were recorded blind --- the original
label was withheld to a separate key --- under Paper 3's own criterion:
a question counts as engaged if some in-window paper investigates its
own objects, quantities, or claimed relationship, not merely its topic.

The result is unambiguous (Table~\ref{tab:adjudication}). Of the $36$
questions the label calls \texttt{not\_addressed}, $34$ were in fact
engaged --- a recovery rate of $94\%$. Under the strictest reading, in
which only unambiguous direct engagement counts and every partial
judgment is discarded, $25/36 = 69\%$ are still recovered. Both
brackets sit around the $80\%$ recovery Paper 3 itself documents for
its object-level Paper 1 import ($2/10 \to 10/10$), which is
independent corroboration that this is the same measurement bound
rather than an artifact of our judge. Re-running the stratum's utility
comparison under the corrected labels, the effect disappears entirely:
$p = 0.70$ lenient, $p = 0.61$ strict, against $p = 0.004$ under the
original labels.

\begin{table}[t]
\centering
\caption{Adjudication of the 62 high-specificity questions under
citation-graph retrieval, blind to the original label. The recovery
rate is the fraction of \texttt{not\_addressed} questions found to have
been engaged.}
\label{tab:adjudication}
\begin{tabular}{lrrr}
\toprule
Engagement bar & Recovery & Utility effect & $p$ \\
\midrule
Original label (Paper 3) & --- & 1.830 vs.\ 1.925 & \textbf{0.004} \\
Lenient (direct or partial) & 34/36 = 94\% & 1.883 vs.\ 1.913 & 0.70 \\
Strict (direct only) & 25/36 = 69\% & 1.875 vs.\ 1.917 & 0.61 \\
\bottomrule
\end{tabular}
\end{table}

\paragraph{What this settles, and what it costs.} The neglect
interpretation is withdrawn: the one surviving piece of it was an
artifact of a label that misses most engagement with specific
questions. What remains from the scaled run is the mechanism ---
planner utility and community attention are driven with opposite sign
by cluster breadth --- and that is a statement about how the two
allocation policies differ structurally, not evidence that the planner
finds neglected science. Establishing the latter needs an outcome
measure that does not degrade with question specificity, which is now a
concrete requirement rather than an open worry.

The measurement itself is the durable result. On this stratum the
\texttt{addressed} label under-counts engagement by a factor of roughly
three to fifteen depending on the bar, which bounds what any
specificity-related finding computed on it can mean --- including Paper
3's own negative controlled effect of specificity on attention
($\beta = -0.16$), which it already flags as partly instrumental. Our
adjudication, its rules, and its per-question verdicts are released so
the bound can be checked and reused rather than re-derived.

\begin{table}[t]
\centering
\caption{Robustness of the utility/attention anti-correlation.
Within-stratum comparison of top-$K$ utility for addressed vs.\
unaddressed questions (seeded permutation test). The pooled effect
($p = 0.004$) does not survive breadth stratification and concentrates
in the stratum where Paper 3's label is a documented lower bound.}
\label{tab:robust}
\begin{tabular}{llrrrr}
\toprule
Stratum & Cell & $n$ & Addressed & Unaddressed & $p$ \\
\midrule
--- & pooled & 374 & 1.887 & 1.919 & \textbf{0.004} \\
\midrule
Cluster breadth & 2-paper & 71 & 1.926 & 1.936 & 0.25 \\
                & 3-paper & 59 & 1.915 & 1.913 & 0.88 \\
                & 4+-paper & 244 & 1.877 & 1.914 & 0.08 \\
\midrule
Specificity & low (0--1) & 249 & 1.910 & 1.919 & 0.11 \\
            & mid (2) & 63 & 1.854 & 1.911 & 0.60 \\
            & high (3+) & 62 & 1.830 & 1.925 & \textbf{0.004} \\
\bottomrule
\end{tabular}
\end{table}

\paragraph{Scope notes.} The distributions are narrower than a mature
benchmark should produce, because 2--4-sentence clusters support only
schematic graphs; full-text evidence extraction would widen the spread.
The conversion rules are heuristics operating on cluster-level flags,
not per-sentence stance; a stance-aware mapping is the other obvious
refinement. Both limitations are properties of the imported data's
granularity, not of the benchmark protocol. One further caution
inherited wholesale from Paper 3: outcome rates are functions of the
(corpus, retriever, judge) triple --- re-judging its ten-question
anchor on a $2.4\times$ larger corpus flipped four labels in both
directions --- so every comparison here is made within one instance and
none of the absolute levels should be read across instances.

\section{Implementation}

The reference implementation is pure Python (standard library plus
PyYAML), organized as ten modules that communicate only through shared
dataclasses and the graph interface; estimators and generators are
injected, never imported across modules. All tunable behavior ---
utility weights, gap thresholds, the action catalog, discriminative
powers, source novelties, resource sets --- lives in YAML or documented
constructor parameters. The test suite (57 unit tests) covers question
loading, the graph, gap analysis, action generation, every utility
component, planner ranking, belief updating, each evaluation metric, and
the end-to-end benchmark, the cluster-to-graph conversion rules, and
the statistical helpers (permutation test, rank correlation,
stratification). An experiment script reproduces Section~\ref{sec:case}
in one command; a second reproduces the full 374-question run of
Section~\ref{sec:scale} including its robustness stratification; and a
third loads Paper 3's released question and outcome-label records
directly into planner objects (their extra fields --- source, cutoff,
provenance cluster --- are preserved in question metadata). Paper 3's
tension subset is vendored verbatim so the scaled benchmark runs
offline and reproduces byte-for-byte.

\section{Limitations and Future Work}

\begin{itemize}
  \item \textbf{Two-outcome information gain.} The preposterior model
  collapses an action's result to favors/disfavors with a per-type
  likelihood ratio. Multi-outcome models, and ratios conditioned on the
  specific hypothesis pair, are natural extensions behind the same
  interface.
  \item \textbf{Type-and-target ground-truth matching.} Resolution
  efficiency currently matches recommendations to hidden evidence by
  action type and target hypothesis; semantic matching would credit
  plans that reach the same evidence by a different route.
  \item \textbf{Exhaustive, exclusive hypothesis sets.} The Bayesian
  updater assumes the hypothesis set partitions the possibility space;
  an explicit ``none of the above'' mass is a small but valuable
  addition.
  \item \textbf{Static candidate space.} Actions are instantiated from
  gaps in the current graph; sequential planning over multi-step
  campaigns (with evidence-dependent branching) is the obvious next
  step, and the offline design makes it benchmarkable the same way.
  \item \textbf{Benchmark granularity.} The scaled run
  (Section~\ref{sec:scale}) covers $374$ questions, but its graphs are
  schematic --- built from 2--4 abstract sentences via cluster-level
  flags, without per-sentence stance --- and its labels lack the action
  types needed for resolution efficiency at scale. The two concrete
  refinements are stance-aware edge mapping and mining action types
  from the papers that historically resolved each question; only the
  fully specified case study currently exercises all seven metrics.
  \item \textbf{Empirical grounding is correlational and one-domain.}
  Paper 3's effects come from astronomy alone, from abstract-level
  retrieval, and from an LLM judge with measured (low) inter-rater
  ceilings ($\kappa = 0.21$ against a second judge, against a
  human--human ceiling of $\kappa = 0.17$ on an analogous taxonomy);
  the calibrated configuration inherits those caveats. Its own
  prospective instance (frozen 2026-08, scored 2027--2030) will test
  the imported prior on a future no model has seen.
  \item \textbf{The outcome label bounds what can be concluded at
  scale.} Our at-scale comparisons run against a retrieval-bounded
  label that under-counts engagement severely and non-uniformly:
  Section~\ref{sec:adjudication} measures $69$--$94\%$ under-counting
  on the high-specificity stratum. Any result computed on this label
  that correlates with question specificity is suspect by default,
  which is why the neglect interpretation is withdrawn rather than
  hedged. Adjudicating the remaining $312$ tension questions would
  extend the bound beyond the one stratum where we measured it.
  \item \textbf{The adjudication is model-judged, on one retrieval
  channel.} Our verdicts come from a language model applying Paper 3's
  criterion, not from human adjudication, so they inherit the judge
  limitation Paper 3 documents ($\kappa = 0.21$ model--model, against a
  $\kappa = 0.17$ human--human ceiling). What makes the comparison
  informative is that the \emph{retrieval} differs while the criterion
  is held fixed, and that the recovery brackets Paper 3's own $80\%$
  human-adjudicated recovery on an independent set. Citation-graph
  retrieval also has its own blind spot, symmetric to the one it fixes:
  it cannot see engagement that never cites the generating papers, so
  our recovery rates are themselves lower bounds.
\end{itemize}

\section{Conclusion}

We presented an offline, domain-independent framework that turns ``what
should we investigate next?'' into a concrete, configurable, and ---
crucially --- historically evaluable planning problem. By separating
evidence representation, gap analysis, action generation, utility
estimation, and belief updating behind small interfaces, the framework
serves as a reference implementation that individual research programs
can adapt by swapping components and editing configuration, not code.
The historical benchmark reframes evaluation from plausibility to
counterfactual usefulness: a good evidence planner is one that would
have gotten science to the answer sooner.

The connection to Paper 3 closes a loop the series has been building
toward. Paper 3 shows that what the community \emph{will} engage is
predictable and community-relational, while what \emph{resolves}
questions is discriminative --- quantified, contradictory, recurrent
evidence. This framework operationalizes the second finding as a
planning objective, imports the first as a calibrated configuration,
and on the series' own canonical case study recommends, first, the
exact action family that history used to settle the question. Run at
scale over all $374$ of Paper 3's tension questions, the planner's
utility turns out to anti-correlate with subsequent community
attention: it prizes precisely the evidence-poor questions the
community left behind. We tested that reading and it failed. Blind
re-adjudication of the stratum carrying it, against a retrieval channel
the original label lacked, recovered $94\%$ of the questions the label
called unaddressed and dissolved the effect entirely. What survives is
narrower and still worth having: planner utility and community
attention are driven with opposite sign by the same structural feature
--- a statement about how two allocation policies differ, not evidence
that either finds better science.

The more useful result is the one we did not set out to produce. A
planner benchmark turns out to be a sensitive instrument for auditing
the outcome labels it is scored against, because it concentrates
disagreement exactly where a label is weakest. The bound we measured
--- a factor of three to fifteen in missed engagement for specific
questions --- constrains any specificity-related claim built on that
label, ours and the source dataset's alike. Whether a
resolution-driven planner genuinely points at neglected science remains
open, and answering it now has a precondition this work made concrete:
an outcome measure that does not degrade with question specificity.
Paper 3's pre-registered prospective instance (frozen 2026, scored
2027--2030) offers the time-graded half of that test; the retrieval
half is a solved problem we have implemented and released.

\section*{Reproducibility}

The framework, configurations, vendored benchmark inputs, and the
scaled-run outputs are released. Every module is a single file with
unit tests; the case study and the 374-question benchmark each
reproduce in one command
(\texttt{experiments/run\_benchmark.py},
\texttt{experiments/paper3\_benchmark.py}). The graph-conversion rules,
the permutation tests, and the stratification cells are implemented in
the released scripts rather than described only in prose. Paper 3's
dataset is used verbatim as published, and is additionally available on
the Hugging Face Hub at
\texttt{huiluckylucky/scientific-question-outcomes}.

\section*{Declaration of generative AI and AI-assisted technologies in
the writing process}

During the preparation of this work the author used Claude (Anthropic)
to draft and edit manuscript prose, to typeset the manuscript in
\LaTeX{}, and to assist in writing the framework and benchmark code in
the released repository. After using this tool, the author reviewed and
edited the content as needed and takes full responsibility for the
content of the publication.

\begin{thebibliography}{9}

\bibitem{paper1}
Hui Mao.
\emph{Evidence-Based Scientific Question Discovery: A Framework with
Historical Backtesting.}
arXiv:2608.09968, 2026. Referred to as ``Paper 1''.

\bibitem{paper2}
Hui Mao.
\emph{Historical Backtesting for Scientific Question Discovery: A
Protocol and Astronomy Pilot.}
arXiv:2608.16795 (v1.1, sealed), 2026. Referred to as ``Paper 2''.

\bibitem{paper3}
Hui Mao.
\emph{What Makes a Scientific Question Succeed? Predicting Future
Attention, Resolution, and Premise Revision from Question Structure and
Evidence Context.}
2026. Referred to as ``Paper 3''.

\bibitem{hindsight}
Bo Jiang.
\emph{HindSight: Evaluating LLM-Generated Research Ideas via Future
Impact.}
arXiv:2603.15164, 2026.

\bibitem{si2025execution}
Chenglei Si, Tatsunori Hashimoto, and Diyi Yang.
\emph{The Ideation--Execution Gap: Execution Outcomes of LLM-Generated
versus Human Research Ideas.}
arXiv:2506.20803, 2025.

\bibitem{si2024novel}
Chenglei Si, Diyi Yang, and Tatsunori Hashimoto.
\emph{Can LLMs Generate Novel Research Ideas? A Large-Scale Human Study
with 100+ NLP Researchers.}
arXiv:2409.04109, 2024.

\bibitem{rzhetsky2015}
Andrey Rzhetsky, Jacob G. Foster, Ian T. Foster, and James A. Evans.
\emph{Choosing experiments to accelerate collective discovery.}
Proceedings of the National Academy of Sciences, 112(47):14569--14574,
2015.

\bibitem{foster2015}
Jacob G. Foster, Andrey Rzhetsky, and James A. Evans.
\emph{Tradition and Innovation in Scientists' Research Strategies.}
American Sociological Review, 80(5):875--908, 2015.

\bibitem{kitano2021}
Hiroaki Kitano.
\emph{Nobel Turing Challenge: creating the engine for scientific
discovery.}
npj Systems Biology and Applications, 7:29, 2021.

\end{thebibliography}

\end{document}
